\section{The \texorpdfstring{$\W5$}{W5} instance as a predictive transform}\label{sec:predictive}

The one-bit state has a direct operational meaning under $\W5$.  Fix a context $a$, and list the
positions at which it occurs as
\[
  k_0<k_1<k_2<\cdots .
\]
Let $y_j=x_{k_j}$ be the symbol that follows that occurrence of $a$.  If the initial cell is
$q_a$, $\cxor_{\W5}$ emits
\begin{equation}\label{eq:w5-per-context}
  r_{k_0}=y_0+q_a,
  \qquad
  r_{k_j}=y_j+y_{j-1}\quad(j\ge1).
\end{equation}
The first visit is predicted from the initial table.  Every later visit compares the current
successor with the successor observed on the preceding visit to the same context.  Thus $\W5$ is
contextual differencing: after the first visit, a zero says exactly that the association
$a\mapsto y$ has repeated.

This gives a more accurate description than saying that $\W5$ detects periodic substrings.
$\W5$ suppresses repeated context-to-successor associations.  Periodic and locally repetitive
sources are important special cases: once an $m$-bit context recurs with a stable successor,
subsequent visits to that context emit zero.  But periodicity is neither necessary nor sufficient
as a source-wide description.  Different contexts may recur at irregular intervals, and a
repeated context whose successor alternates will produce errors rather than zeros.

\subsection{What determines zero enrichment}

The residual depends jointly on the source and the model order.
\begin{itemize}
  \item \emph{Context order.}  A small $m$ merges histories that may have different successors;
        a large $m$ separates them but creates more cells to learn.
  \item \emph{Input length and recurrence.}  There are $2^m$ possible contexts.  On a short input,
        many are never revisited, so their one-bit memories cannot contribute a learned
        prediction.  A context seen only once is predicted solely from its initial cell.
  \item \emph{Successor stability.}  Recurrence helps only when repeated occurrences tend to have
        the same following bit.  Equation~\eqref{eq:w5-per-context} makes this dependence exact.
  \item \emph{Initialisation.}  The initial table affects the first visit to each context.  The
        zero table used by the reference implementation initially predicts zero everywhere;
        $\W5$ overwrites that choice after one observation.
\end{itemize}

The example in the introduction illustrates these conditions rather than proving a general
gain.  Its 33 symbols are long enough to revisit several four-bit contexts, and enough of those
contexts repeat their successor that the residual contains 20 zeros instead of 7.  On another
source, or with another $m$, $\W5$ may leave the marginal balance unchanged or make it less
favourable.  Prediction performance is therefore empirical even though the transform is exact.

\subsection{Why this is not compression}

$\cxor_w$ maps $n$ input bits to $n$ residual bits, and Theorem~\ref{thm:causal-inverse} says that
it permutes the $2^n$ possible strings for every wiring.  It does not reduce length or information
entropy by itself.
The intended architecture, if the residual is useful, is
\[
  x
  \longrightarrow
  \text{$\cxor_{\W5}$ residual}
  \longrightarrow
  \text{statistical model and entropy coder}.
\]
A downstream coder might exploit a residual with a strong zero bias, long zero runs, or other
simple structure more easily than it exploits the original context dependence.  Whether it
actually does so depends on that coder and on the source; it is an experimental question not
answered by the algebra in this paper.

The role is analogous in spirit to other reversible preprocessing transforms.  The
Burrows--Wheeler transform, for example, is valuable because it presents existing dependence in
a form suitable for simple back ends, not because the permutation itself shortens the data
\cite{burrows1994block}.  $\cxor_{\W5}$ uses a causal predictor rather than block sorting, so the analogy
should not be taken further.

Zero frequency, invertibility, and compressed size therefore answer different questions.  The
first measures hard-prediction success for a specific run.  The second is guaranteed by causal
decoding.  The third can be known only after a downstream representation has been specified.
